\subsection{THE PRINCIPLE OF TRANSFINITE INDUCTION}

LEMMA 2. Let E be a well-ordered set and let S be a set of segments of E with the following properties: (1) every union of segments belonging to S belongs to S; (2) if \(S_{x} \in S\) , then \(S_{x} \cup \{x\} \in S\) . Then every segment of E belongs to S.

Suppose that there are segments of E which do not belong to S, and let S be the smallest of them (no. 1, Proposition 2). If S has no greatest element, then S is the union of the segments of S distinct from S itself, and these segments belong to S by virtue of the definition of S; hence \(S \in S\) , which is absurd. If, on the other hand, S has a greatest element a, then \(S = S_{a} \cup \{a\}\) , and since \(S_{a}\) is a segment of S distinct from S, we have \(S_{a} \in S\) ; but then also \(S \in S\) , which again is absurd.

For greater convenience we shall place ourselves in a theory \(\mathfrak{G}\) in which \(\mathbf{E}\) is a set well-ordered by a relation written \(x \leqslant y\) . We have then the following criteria:

C59. (Principle of transfinite induction). Let \(\mathbf{R}\{x\}\) be a relation in \(\mathfrak{G}\) ( \(x\) not being a constant of \(\mathfrak{G}\) ) such that the relation

\[
(x \in \mathrm{E} \text { and } (\forall y) ((y \in \mathrm{E} \text { and } y <   x) \Longrightarrow \mathrm{R} \{\left. y \right\})) \Longrightarrow \mathrm{R} \{x \}
\]

is a theorem in \(\mathfrak{G}\) . Under these conditions the relation \((x\in \mathbf{E})\Rightarrow \mathbb{R}\{x\}\) is a theorem in \(\mathfrak{G}\) .

Let \(\mathfrak{S}\) be the set of segments S of E such that \((y\in S)\Longrightarrow R\{y\}\) . It is clear that every union of segments belonging to \(\mathfrak{S}\) also belongs to \(\mathfrak{S}\) . On the other hand, if \(S_x\in \mathfrak{S}\) , we have \(R\{x\}\) by hypothesis; hence \((y\in S_x\cup \{y\})\Longrightarrow R\{y\}\) by the method of disjunction of cases. Hence (Lemma 2) \(E\in \mathfrak{S}\) , which proves the criterion.

In the applications of C59, the relation

\[
x \in \mathrm{E} \text {   and   } (\forall y) ((y \in \mathrm{E} \text {   and   } y <   x) \Longrightarrow \mathrm{R} \{y \})
\]

is usually called the ``inductive hypothesis''.

In what follows, for every mapping \(g\) of a segment \(S\) of \(E\) into a set \(F\) , and for each \(x \in S\) we shall denote by \(g^{(x)}\) the mapping of the segment \(S_x = ] \leftarrow, x[\) of \(E\) onto \(g(S_x)\) which coincides with \(g\) on \(S_x\) . With this notation we have

C60. (Definition of a mapping by transfinite induction.) Let \(u\) be a letter, \(T\{u\}\) a term in the theory \(\mathfrak{G}\) . There exists a set \(U\) and a mapping \(f\) of \(E\) onto \(U\) such that for all \(x \in E\) we have \(f(x) = T\{f^{(x)}\}\) . Furthermore, the set \(U\) and the mapping \(f\) are uniquely determined by these conditions.

Let us first prove the uniqueness. Suppose that \(f'\) and \(U'\) also satisfy the conditions of the criterion. Let \(\mathfrak{S}\) be the set of segments \(S\) of \(E\) such that \(f\) and \(f'\) coincide on \(S\) . It is clear that every union of segments belonging to \(\mathfrak{S}\) also belongs to \(\mathfrak{S}\) . On the other hand, if \(S_x \in \mathfrak{S}\) , then \(f\) and \(f'\) agree on \(S_x\) and therefore \(f^{(x)} = f'^{(x)}\) ; consequently

\[
f (x) = \mathrm{T} \left\{f ^ {(x)} \right\} = \mathrm{T} \left\{f ^ {\prime (x)} \right\} = f ^ {\prime} (x),
\]

which shows that \(\mathbf{S}_x\cup \{x\} \in \mathfrak{S}\) . It follows that \(\mathbf{E}\in \mathfrak{S}\) (Lemma 2), hence \(f^{\prime} = f\) and \(\mathbf{U}' = f'(\mathbf{E}) = f(\mathbf{E}) = \mathbf{U}\) .

Now let \(\mathfrak{S}_1\) denote the set of segments S of E for which there exists a set \(U_S\) and a mapping \(f_S\) of S onto \(U_S\) such that for all \(x \in S\) we have \(f_S(x) = T \{ f^{(x)} \}\) . For each \(S \in S_1\) , \(f_S\) and \(U_S\) are uniquely determined, by the first part of the proof; in particular, if \(S'\) and \(S''\) are two segments belonging to \(\mathfrak{S}_1\) such that \(S' \subset S''\) , then \(f_S'\) is the mapping of \(S'\) onto \(f_{S'}(S')\) which agrees with \(f_{S'}\) on \(S'\) . From this remark it follows that every union of segments belonging to \(\mathfrak{S}_1\) also belongs to \(\mathfrak{S}_1\) (Chapter II, §4, no. 6, Proposition 7). On the other hand, if \(S_x \in \mathfrak{S}_1\) , we define on \(S = S_x \cup \{x\}\) a function \(f_S\) extending \(f_{S_x}\) by putting

\[
f _ {\mathbf {S}} (x) = \mathrm{T} \left\{f _ {\mathbf {s} _ {x}} \right\}
\]

(Chapter II, § 4, no. 7, Proposition 8); since \(f_{\mathbb{S}}^{(x)} = f_{\mathbb{S}_x}\) , it is obvious that \(\mathbf{S}_x \cup \{x\} \in \mathfrak{S}_1\) . Hence (Lemma 2) \(\mathbf{E} \in \mathfrak{S}_1\) , and the proof is complete.

Usually this criterion is applied in situations where there exists a set F such that for every mapping h of a segment of E onto a subset of F we have \(T\{h\}\in F\) . Then the set U obtained by applying C60 is a subset of F. For, with the notation used above, let \(S_{2}\) be the subset of \(S_{1}\) consisting of segments S of E such that \(U_{S}\subset F\) . It is evident that every union of segments belonging to \(S_{2}\) also belongs to \(S_{2}\) ; on the other hand, the hypothesis on F implies that if \(S_{x}\in S_{2}\) , we have \(S_{x}\cup\{x\}\in S_{2}\) . The assertion now follows from Lemma 2.

